\chapter{2008年陕西卷（理）}

\section{单选题}

\begin{problem}
复数 \(\frac{\i(2 + \i)}{1 - 2\i}\) 等于（\quad）

\choices
    {\(\i\)}
    {\(-\i\)}
    {\(1\)}
    {\(-1\)}
\end{problem}

\begin{answer}
D
\end{answer}

\begin{problem}
已知全集 \(U = \{1, 2, 3, 4, 5\}\)，集合 \(A = \{x \mid x^2 - 3x + 2 = 0\}\)，\(B = \{x \mid x = 2a, a \in A\}\)，则集合 \(\complement_U(A \cup B)\) 中元素的个数为（\quad）

\choices
    {\(1\)}
    {\(2\)}
    {\(3\)}
    {\(4\)}
\end{problem}

\begin{answer}
B
\end{answer}

\begin{problem}
\(\triangle ABC\) 的内角 \(A\)，\(B\)，\(C\) 的对边分别为 \(a\)，\(b\)，\(c\). 若 \(c = \sqrt{2}\)，\(b = \sqrt{6}\)，\(B = 120^{\circ}\)，则 \(a\) 等于（\quad）

\choices
    {\(\sqrt{6}\)}
    {\(2\)}
    {\(\sqrt{3}\)}
    {\(\sqrt{2}\)}
\end{problem}

\begin{answer}
D
\end{answer}

\begin{problem}
已知 \(\{a_{n}\}\) 是等差数列，\(a_{1} + a_{2} = 4\)，\(a_{7} + a_{8} = 28\)，则该数列前 \(10\) 项和 \(S_{10}\) 等于（\quad）

\choices
    {\(64\)}
    {\(100\)}
    {\(110\)}
    {\(120\)}
\end{problem}

\begin{answer}
B
\end{answer}

\begin{problem}
直线 \(\sqrt{3} x - y + m = 0\) 与圆 \(x^{2} + y^{2} - 2x - 2 = 0\) 相切，则实数 \(m\) 等于（\quad）

\choices
    {\(\sqrt{3}\) 或 \(-\sqrt{3}\)}
    {\(-\sqrt{3}\) 或 \(3\sqrt{3}\)}
    {\(-3\sqrt{3}\) 或 \(\sqrt{3}\)}
    {\(-3\sqrt{3}\) 或 \(3\sqrt{3}\)}
\end{problem}

\begin{answer}
C
\end{answer}

\begin{problem}
“\(a = \frac{1}{8}\)”是“对任意的正数 \(x\)，\(2x + \frac{a}{x} \geqslant 1\)”的（\quad）

\choices
    {充分不必要条件}
    {必要不充分条件}
    {充要条件}
    {既不充分也不必要条件}
\end{problem}

\begin{answer}
A
\end{answer}

\begin{problem}
已知函数 \( f(x) = 2^{x + 3} \)，\( f^{-1}(x) \) 是 \( f(x) \) 的反函数，若 \( mn = 16\)（\(m\)，\(n \in \R^+\)），则 \( f^{-1}(m) + f^{-1}(n) \) 的值为（\quad）

\choices
    {\(-2\)}
    {\(1\)}
    {\(4\)}
    {\(10\)}
\end{problem}

\begin{answer}
A
\end{answer}

\begin{problem}
双曲线 \(\frac{x^2}{a^2} - \frac{y^2}{b^2} = 1\)（\(a > 0\)，\(b > 0\)）的左、右焦点分别是 \(F_1\)，\(F_2\)，过 \(F_1\) 作倾斜角为 \(30^\circ\) 的直线交双曲线右支于 \(M\) 点. 若 \(MF_{2}\) 垂直于 \(x\) 轴，则双曲线的离心率为（\quad）

\choices
    {\(\sqrt{6}\)}
    {\(\sqrt{3}\)}
    {\(\sqrt{2}\)}
    {\(\frac{\sqrt{3}}{3}\)}
\end{problem}

\begin{answer}
B
\end{answer}

\begin{problem}
如图，\(\alpha \perp \beta\)，\(\alpha \cap \beta = l\)，\(A \in \alpha\)，\(B \in \beta\)，\(A\)，\(B\) 到 \(l\) 的距离分别是 \(a\) 和 \(b\)，\(AB\) 与 \(\alpha\)，\(\beta\) 所成的角分别是 \(\theta\) 和 \(\varphi\)，\(AB\) 在 \(\alpha\)，\(\beta\) 内的射影分别是 \(m\) 和 \(n\)，若 \(a > b\)，则（\quad）

\FigureLayoutDeclare{img/2008/shaanxi_science/q9_fig1.png}{5.2cm}{4bp 6bp 5bp 6bp}
\bitmapinlinefigure[0.34\linewidth][max width=0.34\linewidth]{img/2008/shaanxi_science/q9_fig1.png}

\choices
    {\(\theta > \varphi\)，\(m > n\)}
    {\(\theta >\varphi\)，\(m < n\)}
    {\(\theta < \varphi\)，\(m < n\)}
    {\(\theta < \varphi\)，\(m > n\)}
\end{problem}

\begin{answer}
D
\end{answer}

\begin{problem}
已知实数 \(x\)，\(y\) 满足 \(\begin{cases}
y \geqslant 1,\\
y \leqslant 2x - 1,\\
x + y \leqslant m,
\end{cases}\) 如果目标函数 \(z = x - y\) 的最小值为 \(-1\)，则实数 \(m\) 等于（\quad）

\choices
    {\(7\)}
    {\(5\)}
    {\(4\)}
    {\(3\)}
\end{problem}

\begin{answer}
B
\end{answer}

\begin{problem}
定义在 \(\R\) 上的函数 \(f(x)\) 满足 \(f(x + y) = f(x) + f(y) + 2xy\)（\(x\)，\(y \in \R\)），\(f(1) = 2\)，则 \(f(-3)\) 等于（\quad）

\choices
    {\(2\)}
    {\(3\)}
    {\(6\)}
    {\(9\)}
\end{problem}

\begin{answer}
C
\end{answer}

\begin{problem}
为提高信息在传输中的抗干扰能力，通常在原信息中按一定规则加入相关数据组成传输信息. 设定原信息为 \(a_0a_1a_2\)，\(a_i \in \{0, 1\}\)（\(i = 0\)，\(1\)，\(2\)），传输信息为 \(h_0a_0a_1a_2h_1\)，其中 \(h_0 = a_0 \oplus a_1\)，\(h_1 = h_0 \oplus a_2\)，\(\oplus\) 运算规则为：\(0 \oplus 0 = 0\)，\(0 \oplus 1 = 1\)，\(1 \oplus 0 = 1\)，\(1 \oplus 1 = 0\). 例如原信息为 \(111\)，则传输信息为 \(01111\). 传输信息在传输过程中受到干扰可能导致接收信息出错，则下列接收信息一定有误的是（\quad）

\choices
    {\(11010\)}
    {\(01100\)}
    {\(10111\)}
    {\(00011\)}
\end{problem}

\begin{answer}
C
\end{answer}

\section{填空题}

\begin{problem}
\(\displaystyle\lim_{n\to \infty}\frac{(1 + a)n + 1}{n + a} = 2\)，则 \(a =\)\fillinblank{}.
\end{problem}

\begin{answer}
1
\end{answer}

\begin{problem}
长方体 \(ABCD - A_{1}B_{1}C_{1}D_{1}\) 的各顶点都在球 \(O\) 的球面上，其中 \(AB:AD:AA_{1} = 1:1:\sqrt{2}\)，\(A\)，\(B\) 两点的球面距离记为 \(m\)，\(A\)，\(D_{1}\) 两点的球面距离记为 \(n\)，则 \(\frac{m}{n}\) 的值为\fillinblank{}.
\end{problem}

\begin{answer}
\(\dfrac{1}{2}\)
\end{answer}

\begin{problem}
关于平面向量 \(\bs{a}\)，\(\bs{b}\)，\(\bs{c}\)，有下列三个命题：

\begin{circlelist}[itemsep=0pt,topsep=0pt,parsep=0pt,partopsep=0pt]
\item 若 \(\bs{a} \cdot \bs{b} = \bs{a} \cdot \bs{c}\)，则 \(\bs{b} = \bs{c}\)；
\item 若 \(\bs{a} = (1, k)\)，\(\bs{b} = (-2, 6)\)，\(\bs{a} \parallel \bs{b}\)，则 \(k = -3\)；
\item 非零向量 \(\bs{a}\) 和 \(\bs{b}\) 满足 \( |\bs{a}| = |\bs{b}| = |\bs{a} - \bs{b}| \)，则 \(\bs{a}\) 与 \( \bs{a} + \bs{b} \) 的夹角为 \( 60^{\circ} \).
\end{circlelist}

其中真命题的序号为\fillinblank{}.（写出所有真命题的序号）.
\end{problem}

\begin{answer}
\(\circled{2}\)
\end{answer}

\begin{problem}
某地奥运火炬接力传递路线共分\(6\)段. 传递活动分别由\(6\)名火炬手完成. 如果第一棒火炬手只能从甲、乙、丙三人中产生，最后一棒火炬手只能从甲、乙两人中产生，则不同的传递方案共有\fillinblank{}种. （用数字作答）
\end{problem}

\begin{answer}
96
\end{answer}

\section{解答题}

\begin{problem}
已知函数 \( f(x) = 2\sin \frac{x}{4}\cos \frac{x}{4} - 2\sqrt{3}\sin^2 \frac{x}{4} + \sqrt{3} \).
\begin{enumerate}
\item 求函数 \( f(x) \) 的最小正周期及最值；
\item 令 \( g(x) = f\left(x + \frac{\pi}{3}\right) \)，判断函数 \( g(x) \) 的奇偶性，并说明理由.
\end{enumerate}
\end{problem}

\begin{solution}
\begin{enumerate}
\item 利用降幂公式，得
\[
f(x)=\sin\frac{x}{2}-\sqrt{3}\left(1-\cos\frac{x}{2}\right)+\sqrt{3}=\sin\frac{x}{2}+\sqrt{3}\cos\frac{x}{2}=2\sin\left(\frac{x}{2}+\frac{\pi}{3}\right).
\]
因此，函数 \(f(x)\) 的最小正周期为 \(4\pi\)，最大值为 \(2\)，最小值为 \(-2\).

\item 由上式，
\[
g(x)=2\sin\left(\frac{x}{2}+\frac{\pi}{6}+\frac{\pi}{3}\right)=2\sin\left(\frac{x}{2}+\frac{\pi}{2}\right)=2\cos\frac{x}{2}.
\]
所以 \(g(-x)=2\cos\left(-\frac{x}{2}\right)=2\cos\frac{x}{2}=g(x)\)，故 \(g(x)\) 是偶函数.
\end{enumerate}
\end{solution}

\begin{problem}
某射击测试规则为：每人最多射击 \(3\) 次，击中目标即终止射击，第 \(i\) 次击中目标得 \(4-i\) 分（\(i = 1\)，\(2\)，\(3\)），\(3\) 次均未击中目标得 \(0\) 分. 已知某射手每次击中目标的概率为 \(0.8\)，其各次射击结果互不影响.
\begin{enumerate}
\item 求该射手恰好射击两次的概率；
\item 该射手的得分记为 \(\xi\)，求随机变量 \(\xi\) 的分布列及数学期望.
\end{enumerate}
\end{problem}

\begin{solution}
\begin{enumerate}
\item 恰好射击两次，表示第一次未击中、第二次击中，因此所求概率为 \(P=0.2\times0.8=0.16\).

\item 得分为 \(3\) 分、\(2\) 分、\(1\) 分、\(0\) 分的事件分别为第一次击中、第一次未击中而第二次击中、前两次未击中而第三次击中、三次均未击中. 因此 \(P(\xi=3)=0.8\)，\(P(\xi=2)=0.2\times0.8=0.16\)，\(P(\xi=1)=0.2^2\times0.8=0.032\)，\(P(\xi=0)=0.2^3=0.008\).
所以 \(\xi\) 的分布列为：\(P(\xi=0)=0.008\)，\(P(\xi=1)=0.032\)；

\(P(\xi=2)=0.16\)，\(P(\xi=3)=0.8\). 从而
\[
E\xi=0\times0.008+1\times0.032+2\times0.16+3\times0.8=2.752.
\]
\end{enumerate}
\end{solution}

\begin{problem}
三棱锥被平行于底面 \(ABC\) 的平面所截得的几何体如图所示，截面为 \(A_{1}B_{1}C_{1}\)，\(\angle BAC = 90^{\circ}\)，\(A_{1}A \perp\) 平面 \(ABC\)，\(A_{1}A = \sqrt{3}\)，\(AB = \sqrt{2}\)，\(AC = 2\)，\(A_{1}C_{1} = 1\)，\(\frac{BD}{DC} = \frac{1}{2}\).

\begin{enumerate}
\item 证明：平面 \(A_{1}AD\) \(\perp\) 平面 \(BCC_{1}B_{1}\)；
\item 求二面角 \(A-CC_{1}-B\) 的大小.
\end{enumerate}

\bitmapfigure[max width=0.34\linewidth]{img/2008/shaanxi_science/q19_fig1.png}
\end{problem}

\begin{solution}
建立空间直角坐标系，使 \(A\) 为原点，\(AB\)，\(AC\)，\(AA_{1}\) 分别沿 \(x\)，\(y\)，\(z\) 轴正方向. 由已知条件，
\(A=(0,0,0)\)，\(B=(\sqrt{2},0,0)\)，\(C=(0,2,0)\)，\(A_{1}=(0,0,\sqrt{3})\).
由于截面 \(A_{1}B_{1}C_{1}\) 平行于底面 \(ABC\)，且 \(A_{1}C_{1}=1\)，\(AC=2\)，所以截面与底面的相似比为 \(\frac{1}{2}\)，从而
\(B_{1}=\left(\frac{\sqrt{2}}{2},0,\sqrt{3}\right)\)，\(C_{1}=(0,1,\sqrt{3})\).
又因为 \(\frac{BD}{DC}=\frac{1}{2}\)，所以 \(\frac{BD}{BC}=\frac{1}{3}\)，于是
\[
D=B+\frac{1}{3}\overrightarrow{BC}=\left(\frac{2\sqrt{2}}{3},\frac{2}{3},0\right).
\]

\begin{enumerate}
\item 平面 \(A_{1}AD\) 的两个方向向量可取 \(\overrightarrow{AA_{1}}=(0,0,\sqrt{3})\) 和 \(\overrightarrow{AD}=\left(\frac{2\sqrt{2}}{3},\frac{2}{3},0\right)\)，故其法向量可取
\[
\bs n_{1}=(\sqrt{3},-\sqrt{6},0).
\]
平面 \(BCC_{1}B_{1}\) 的两个方向向量可取 \(\overrightarrow{BC}=(-\sqrt{2},2,0)\) 和 \(\overrightarrow{BB_{1}}=\left(-\frac{\sqrt{2}}{2},0,\sqrt{3}\right)\)，故其法向量可取
\[
\bs n_{2}=(2\sqrt{3},\sqrt{6},\sqrt{2}).
\]
于是 \(\bs n_{1}\cdot\bs n_{2}=6-6=0\)，两平面的法向量互相垂直，因此平面 \(A_{1}AD\perp\) 平面 \(BCC_{1}B_{1}\).

\item 二面角 \(A-CC_{1}-B\) 是平面 \(ACC_{1}\) 与平面 \(BCC_{1}B_{1}\) 所成的角. 平面 \(ACC_{1}\) 的法向量可取
\[
\bs n_{3}=\overrightarrow{AC}\times\overrightarrow{AC_{1}}=(2\sqrt{3},0,0).
\]
设二面角的大小为 \(\omega\)，则
\[
\cos\omega=\frac{|\bs n_{2}\cdot\bs n_{3}|}{|\bs n_{2}|\,|\bs n_{3}|}=\frac{12}{\sqrt{20}\cdot2\sqrt{3}}=\frac{\sqrt{15}}{5}.
\]
所以二面角 \(A-CC_{1}-B\) 的大小为 \(\arccos\frac{\sqrt{15}}{5}\).
\end{enumerate}
\end{solution}

\begin{problem}
已知抛物线 \(C: y = 2x^2\)，直线 \(y = kx + 2\) 交 \(C\) 于 \(A\)，\(B\) 两点. \(M\) 是线段 \(AB\) 的中点，过 \(M\) 作 \(x\) 轴的垂线交 \(C\) 于点 \(N\).
\begin{enumerate}
\item 证明：抛物线 \(C\) 在点 \(N\) 处的切线与 \(AB\) 平行；
\item 是否存在实数 \(k\) 使 \(\overrightarrow{NA} \cdot \overrightarrow{NB} = 0\). 若存在，求 \(k\) 的值；若不存在，说明理由.
\end{enumerate}
\end{problem}

\begin{solution}
设直线与抛物线的两个交点的横坐标分别为 \(x_{1}\)，\(x_{2}\)，则 \(x_{1}\)，\(x_{2}\) 是方程 \(2x^{2}-kx-2=0\) 的两个根. 由韦达定理，
\(x_{1}+x_{2}=\frac{k}{2}\)，\(x_{1}x_{2}=-1\).
记 \(s=x_{1}+x_{2}=\frac{k}{2}\)，则
\[
M=\left(\frac{x_{1}+x_{2}}{2},\frac{2x_{1}^{2}+2x_{2}^{2}}{2}\right)=\left(\frac{k}{4},\frac{k^{2}}{4}+2\right).
\]
所以
\[
N=\left(\frac{k}{4},2\left(\frac{k}{4}\right)^{2}\right)=\left(\frac{k}{4},\frac{k^{2}}{8}\right).
\]

\begin{enumerate}
\item 抛物线 \(y=2x^{2}\) 在横坐标为 \(x\) 处的切线斜率为 \(4x\). 因此抛物线在点 \(N\) 处的切线斜率为 \(4\cdot\frac{k}{4}=k\)，而直线 \(AB\) 的斜率也是 \(k\)，故该切线与 \(AB\) 平行.

\item 记 \(d=x_{1}-x_{2}\). 由于方程 \(2x^{2}-kx-2=0\) 的判别式为 \(k^{2}+16>0\)，所以 \(d\neq0\). 又有 \(N=\left(\frac{s}{2},\frac{s^{2}}{2}\right)\)，且
\(x_{1}=\frac{s+d}{2}\)，\(x_{2}=\frac{s-d}{2}\).
从而
\[
\overrightarrow{NA}=\left(\frac{d}{2},\frac{d(2s+d)}{2}\right).
\]
\[
\overrightarrow{NB}=\left(-\frac{d}{2},\frac{d(d-2s)}{2}\right).
\]
于是
\[
\overrightarrow{NA}\cdot\overrightarrow{NB}=\frac{d^{2}}{4}\left(d^{2}-4s^{2}-1\right).
\]
由 \(x_{1}x_{2}=-1\)，得 \(d^{2}=(x_{1}+x_{2})^{2}-4x_{1}x_{2}=s^{2}+4\)，所以
\[
\overrightarrow{NA}\cdot\overrightarrow{NB}=\frac{3d^{2}}{4}(1-s^{2}).
\]
因此 \(\overrightarrow{NA}\cdot\overrightarrow{NB}=0\) 等价于 \(s^{2}=1\)，即 \(\left(\frac{k}{2}\right)^{2}=1\). 所以存在这样的实数，且 \(k=\pm2\).
\end{enumerate}
\end{solution}

\begin{problem}
已知函数 \( f(x) = \frac{kx + 1}{x^2 + c}\)（\(c > 0\)且 \(c \neq 1\)，\(k \in \R\)）恰有一个极大值点和一个极小值点. 其中一个是 \(x = -c\).
\begin{enumerate}
\item 求函数 \( f(x) \) 的另一个极值点；
\item 求函数 \( f(x) \) 的极大值 \(M\) 和极小值 \( m \)，并求 \( M - m \geqslant 1 \) 时 \( k \) 的取值范围.
\end{enumerate}
\end{problem}

\begin{solution}
函数的定义域为 \(\R\)，且
\[
f'(x)=\frac{-kx^{2}-2x+kc}{(x^{2}+c)^{2}}.
\]
由于 \(x=-c\) 是极值点，必有 \(f'(-c)=0\)，即
\[
c\left(2+k-kc\right)=0.
\]
由 \(c>0\)，得 \(k=\frac{2}{c-1}\). 此时
\[
-kx^{2}-2x+kc=-k(x+c)(x-1),
\]
所以函数的另一个极值点为 \(x=1\).

当 \(c>1\) 时，\(k>0\)，由导数符号可知 \(x=-c\) 是极小值点，\(x=1\) 是极大值点；当 \(0<c<1\) 时，\(k<0\)，由导数符号可知 \(x=-c\) 是极大值点，\(x=1\) 是极小值点. 直接计算得
\(f(-c)=-\frac{1}{c(c-1)}\)，\(f(1)=\frac{1}{c-1}\).

\begin{enumerate}
\item 函数的另一个极值点的横坐标为 \(x=1\)，对应的极值点为 \(\left(1,\frac{1}{c-1}\right)\).

\item 当 \(c>1\) 时，
\(M=\frac{1}{c-1}\)，\(m=-\frac{1}{c(c-1)}\)，\(M-m=\frac{c+1}{c(c-1)}\).
此时 \(M-m\geqslant1\) 等价于 \(c^{2}-2c-1\leqslant0\). 结合 \(c>1\)，得 \(1<c\leqslant1+\sqrt{2}\).

当 \(0<c<1\) 时，
\(M=\frac{1}{c(1-c)}\)，\(m=-\frac{1}{1-c}\)，\(M-m=\frac{c+1}{c(1-c)}>1\).
所以此时所有 \(0<c<1\) 均满足条件. 由 \(k=\frac{2}{c-1}\)，当 \(0<c<1\) 时 \(k<-2\)，当 \(1<c\leqslant1+\sqrt{2}\) 时 \(k\geqslant\sqrt{2}\). 因此 \(k\) 的取值范围为
\[
(-\infty,-2)\cup[\sqrt{2},+\infty).
\]
\end{enumerate}
\end{solution}

\begin{problem}
已知数列 \(\{a_{n}\}\) 的首项 \(a_1 = \frac{3}{5}\)，\(a_{n+1} = \frac{3a_n}{2a_n + 1}\)，\(n = 1\)，\(2\)，\(\cdots\).
\begin{enumerate}
\item 求 \(\{a_{n}\}\) 的通项公式；
\item 证明：对任意的 \(x > 0\)，\(a_{n} \geqslant \frac{1}{1 + x} - \frac{1}{(1 + x)^{2}} \left( \frac{2}{3^{n}} - x \right)\)，\(n = 1\)，\(2\)，\(\cdots\)；
\item 证明：\(a_{1} + a_{2} + \dots + a_{n} > \frac{n^{2}}{n + 1}\).
\end{enumerate}
\end{problem}

\begin{solution}
由 \(a_{1}>0\) 及递推关系可知各项均为正数. 令 \(b_{n}=\frac{1}{a_{n}}\)，则 \(b_{1}=\frac{5}{3}\)，\(b_{n+1}=\frac{2a_{n}+1}{3a_{n}}=\frac{2}{3}+\frac{1}{3}b_{n}\).
因此 \(b_{n+1}-1=\frac{1}{3}(b_{n}-1)\)，且 \(b_{1}-1=\frac{2}{3}\).
从而 \(b_{n}-1=\frac{2}{3^{n}}\)，故
\[
a_{n}=\frac{1}{b_{n}}=\frac{1}{1+\frac{2}{3^{n}}}=\frac{3^{n}}{3^{n}+2}.
\]

\begin{enumerate}
\item \(\displaystyle a_{n}=\frac{3^{n}}{3^{n}+2}\).

\item 记 \(t=\frac{2}{3^{n}}\). 对任意 \(x>0\)，有
\[
a_{n}-\left[\frac{1}{1+x}-\frac{1}{(1+x)^{2}}\left(\frac{2}{3^{n}}-x\right)\right]
=\frac{1}{1+t}-\frac{1+2x-t}{(1+x)^{2}}=\frac{(x-t)^{2}}{(1+t)(1+x)^{2}}\geqslant0.
\]
所以
\[
a_{n}\geqslant\frac{1}{1+x}-\frac{1}{(1+x)^{2}}\left(\frac{2}{3^{n}}-x\right).
\]

\item 对上面的不等式从 \(i=1\) 到 \(i=n\) 求和，并令 \(x=\frac{1}{n}>0\). 由
\[
2\sum_{i=1}^{n}\frac{1}{3^{i}}=1-\frac{1}{3^{n}},
\]
可得
\[
a_{1}+a_{2}+\cdots+a_{n}\geqslant\frac{n}{1+x}-\frac{1}{(1+x)^{2}}\left(1-\frac{1}{3^{n}}-nx\right)=\frac{n+2nx-1+3^{-n}}{(1+x)^{2}}.
\]
代入 \(x=\frac{1}{n}\)，得
\[
 a_{1}+a_{2}+\cdots+a_{n}\geqslant\frac{n^{2}}{(n+1)^{2}}\left(n+1+\frac{1}{3^{n}}\right)=\frac{n^{2}}{n+1}+\frac{n^{2}}{3^{n}(n+1)^{2}}>\frac{n^{2}}{n+1}.
\]
结论得证.
\end{enumerate}
\end{solution}
